\documentclass[aspectratio=169]{beamer}
\usepackage{tikz}
\usepackage{graphicx}
\usepackage{listings}
\usepackage{xcolor}
\usepackage{hyperref}
\usepackage[spanish,es-noquoting]{babel}
\usepackage[round]{natbib}

\lstset{
  basicstyle=\ttfamily\small,
  keywordstyle=\color{blue},
  stringstyle=\color{red},
  commentstyle=\color{gray},
  showstringspaces=false
}

% Clean Theme Setup & Optimization
\usetheme{Warsaw}
\usecolortheme{default}
\usefonttheme{serif}

% Remove top headline bar and navigation clutter to maximize slide text space
\setbeamertemplate{headline}{}
\setbeamertemplate{navigation symbols}{}

% Slim, clean footline bar
\setbeamertemplate{footline}{
  \leavevmode%
  \hbox{%
    \begin{beamercolorbox}[wd=.333333\paperwidth,ht=2.25ex,dp=1ex,left]{author in head/foot}%
      \hspace*{2ex}\insertshortauthor
    \end{beamercolorbox}%
    \begin{beamercolorbox}[wd=.333333\paperwidth,ht=2.25ex,dp=1ex,center]{title in head/foot}%
      \insertshorttitle
    \end{beamercolorbox}%
    \begin{beamercolorbox}[wd=.333333\paperwidth,ht=2.25ex,dp=1ex,right]{date in head/foot}%
      \insertframenumber{} / \inserttotalframenumber\hspace*{2ex}
    \end{beamercolorbox}%
  }%
  \vskip0pt%
}

% Font size and list spacing optimizations to prevent text overflow
\setbeamertemplate{itemize/enumerate body begin}{\small}
\setbeamertemplate{itemize/enumerate subbody begin}{\footnotesize}
\setbeamertemplate{itemize/enumerate subsubbody begin}{\scriptsize}

\makeatletter
\g@addto@macro\itemize{%
  \setlength{\itemsep}{2pt}%
  \setlength{\parsep}{1pt}%
  \setlength{\topsep}{2pt}%
  \setlength{\partopsep}{0pt}%
}
\makeatother

\graphicspath{{../}{./}}
\setbeamertemplate{caption}[numbered]

% Title info
\title[Análisis del Territorio y Aten]{Análisis del Territorio y Atención Primaria en Salud y Fisioterapia}
\subtitle{Optimizando la Equidad y el Acceso a los Servicios}
\author{Dr. [Su Nombre]}
\institute{PhD en Salud Pública}
\date{\today}



\begin{document}

% Title slide
\begin{frame}
  \titlepage
\end{frame}

% Reset background for normal slides
\usebackgroundtemplate{}


\begin{frame}[allowframebreaks,fragile]{Análisis del Territorio y Atención Primaria en Salud y Fisioterapia}

\textbf{Optimizando la Equidad y el Acceso a los Servicios}

Dr. [Su Nombre], PhD en Salud Pública

Mayo 2024

\end{frame}

\begin{frame}[allowframebreaks,fragile]{Objetivos de la Sesión y Contexto General}


\textbf{Objetivos}
\begin{itemize}
  \item \textbf{Comprender} la relevancia del análisis territorial en la Atención Primaria de Salud (APS) y la Fisioterapia.
  \item \textbf{Identificar} cómo el análisis espacial mejora la distribución y efectividad de los servicios de salud.
  \item \textbf{Reconocer} el rol fundamental del fisioterapeuta en la APS y su impacto en la equidad en salud.
\end{itemize}
\textbf{Contexto General}
\begin{itemize}
  \item El análisis territorial es crucial para entender las dinámicas de salud-enfermedad en una población.
  \item Permite optimizar la provisión de servicios y garantizar la \textbf{equidad en el acceso}.
  \item Desde Salud Pública, identifica disparidades geográficas, socioeconómicas y culturales.
  \item La integración efectiva de la fisioterapia en la APS es fundamental para la prevención, promoción y manejo de condiciones.
\end{itemize}
\end{frame}

\begin{frame}[allowframebreaks,fragile]{Marco Conceptual: Análisis Territorial en Salud Pública}


\textbf{¿Qué es el Análisis Territorial en Salud?}
\begin{itemize}
  \item \textbf{Definición:} Estudio de las dinámicas de salud-enfermedad, recursos sanitarios y factores determinantes dentro de un área geográfica específica.
  \item \textbf{Propósito:} Identificar patrones de distribución de enfermedades, desigualdades en salud y barreras de acceso.
\end{itemize}
\textbf{Metodologías Clave}
\begin{itemize}
  \item \textbf{Sistemas de Información Geográfica (SIG):} Herramientas para mapear, analizar y visualizar datos con componentes espaciales.
  \item \textbf{Mapeo de recursos:} Localización de centros de salud, personal sanitario y equipamiento.
  \item \textbf{Análisis demográfico y socioeconómico:} Caracterización de la población y sus vulnerabilidades.
\end{itemize}
\textbf{Importancia}
\begin{itemize}
  \item \textbf{Informar políticas:} Diseño de estrategias de salud pública basadas en evidencia local.
  \item \textbf{Asignación de recursos:} Distribución eficiente de personal y equipamiento.
  \item \textbf{Diseño de intervenciones:} Adaptación de programas a las necesidades y contextos específicos de cada comunidad.
\end{itemize}
\end{frame}

\begin{frame}[allowframebreaks,fragile]{Marco Conceptual: Atención Primaria de Salud (APS)}


\textbf{Fundamentos de la APS}
\begin{itemize}
  \item \textbf{Definición:} El primer nivel de contacto de los individuos, la familia y la comunidad con el sistema nacional de salud.
  \item \textbf{Principios:}
  \begin{itemize}
    \item \textbf{Accesibilidad:} Proximidad y facilidad de acceso a los servicios.
    \item \textbf{Integralidad:} Abordaje holístico de la salud, cubriendo prevención, promoción, tratamiento y rehabilitación.
    \item \textbf{Continuidad:} Coordinación entre los diferentes niveles de atención.
    \item \textbf{Participación comunitaria:} Involucración activa de la población en la gestión de su salud.
  \end{itemize}
\end{itemize}
\textbf{El Rol de la Fisioterapia en la APS}
\begin{itemize}
  \item \textbf{Intervención temprana:} Abordaje de condiciones musculoesqueléticas y otras disfunciones antes de que se agraven.
  \item \textbf{Prevención:} Programas para prevenir lesiones, caídas y enfermedades crónicas.
  \item \textbf{Promoción de la salud:} Educación sobre estilos de vida saludables, actividad física y ergonomía.
  \item \textbf{Rehabilitación:} Apoyo en la recuperación funcional y mejora de la calidad de vida.
\end{itemize}
\end{frame}

\begin{frame}[allowframebreaks,fragile]{Marco Conceptual: Fisioterapia Basada en el Territorio}


\textbf{Enfoque Adaptativo de la Fisioterapia}
\begin{itemize}
  \item \textbf{Concepto:} Diseño e implementación de servicios de fisioterapia que responden directamente a las necesidades, características y barrates específicas de una comunidad local.
\end{itemize}
\textbf{Consideraciones Territoriales Clave}
\begin{itemize}
  \item \textbf{Factores Geográficos:}
  \begin{itemize}
    \item Distancia a los centros de salud.
    \item Disponibilidad de transporte.
    \item Accesibilidad física de las instalaciones.
  \end{itemize}
  \item \textbf{Factores Socioeconómicos:}
  \begin{itemize}
    \item Nivel de ingresos y capacidad de pago.
    \item Barreras culturales y lingüísticas.
    \item Nivel educativo y alfabetización en salud.
  \end{itemize}
  \item \textbf{Factores Culturales:}
  \begin{itemize}
    \item Creencias y prácticas de salud locales.
    \item Estigmatización de ciertas condiciones o grupos poblacionales.
    \item Necesidad de servicios culturalmente competentes.
  \end{itemize}
\end{itemize}
\textbf{Beneficios}
\begin{itemize}
  \item \textbf{Mayor equidad:} Reducción de disparidades en el acceso y los resultados de salud.
  \item \textbf{Intervenciones pertinentes:} Adaptación que aumenta la adherencia y efectividad del tratamiento.
  \item \textbf{Optimización de recursos:} Direccionamiento de la fisioterapia donde más se necesita.
\end{itemize}
\end{frame}

\begin{frame}[allowframebreaks,fragile]{Casos de Estudio y Ejemplos Prácticos Aplicados}


\textbf{1. Análisis Espacial para Optimizar Recursos en APS}
\begin{itemize}
  \item \textbf{Referencia:} Nascimento, M., \& Rodrigues-Júnior, A. (2014). \textit{Análisis espacial del dengue y la atención primaria de salud en Alfenas, Minas Gerais, Brasil.}
  \begin{itemize}
    \item \textbf{Hallazgo clave:} Demostró cómo el análisis espacial identifica áreas con mayor carga de enfermedad (dengue) y las correlaciona con la calidad o el acceso a los servicios de APS.
  \end{itemize}
  \item \textbf{Aplicación en Fisioterapia:}
  \begin{itemize}
    \item Uso de SIG para mapear la distribución de fisioterapeutas vs. la prevalencia de condiciones musculoesqueléticas por barrio o región.
    \item Identificación de ``zonas calientes'' (hotspots) con alta necesidad y baja oferta de servicios de fisioterapia.
  \end{itemize}
  \item \textbf{Impacto:} Focalizar la apertura de nuevas clínicas, programas móviles o la asignación de personal donde el impacto será mayor, mejorando el acceso y la eficiencia.
\end{itemize}
\end{frame}

\begin{frame}[allowframebreaks,fragile]{Casos de Estudio y Ejemplos Prácticos Aplicados}


\textbf{2. Fisioterapia para Poblaciones Vulnerables desde la APS}
\begin{itemize}
  \item \textbf{Referencia:} Zamarreño, F. (2025). \textit{Salud, territorio y mujeres trans: miradas desde la Atención Primaria.}
  \begin{itemize}
    \item \textbf{Hallazgo clave:} Las barreras territoriales (falta de servicios culturalmente competentes, estigma social en ciertas áreas) impactan significativamente el acceso y la equidad en salud para grupos marginados.
  \end{itemize}
  \item \textbf{Aplicación en Fisioterapia:}
  \begin{itemize}
    \item Desarrollar servicios de fisioterapia itinerantes o comunitarios en áreas remotas o con poblaciones de difícil acceso.
    \item Capacitar a fisioterapeutas en competencia cultural para atender las necesidades específicas de comunidades indígenas, migrantes o LGBTQ+.
    \item Creación de espacios seguros y libres de estigma dentro de la APS para todas las identidades.
  \end{itemize}
  \item \textbf{Impacto:} Reducción de disparidades en salud, mejora de la experiencia del paciente y promoción de la inclusión social.
\end{itemize}
\end{frame}

\begin{frame}[allowframebreaks,fragile]{Ejercicios Prácticos y Actividades Interactivas}


\textbf{Ejercicio 1: Mapeo de Necesidades Comunitarias (10 minutos)}
\begin{itemize}
  \item \textbf{Formen grupos} de 3-4 personas.
  \item \textbf{Elijan} una comunidad local (real o hipotética) que conozcan o imaginen.
  \item \textbf{Identifiquen} 3-5 factores territoriales que podrían influir en el acceso o la calidad de los servicios de fisioterapia en esa comunidad.
  \begin{itemize}
    \item \textit{Ejemplos:} Falta de transporte público, alta población de adultos mayores, presencia de barreras físicas, estigma hacia ciertas condiciones, etc.
  \end{itemize}
  \item \textbf{Discutan:} ¿Cuál de estos factores consideran el más crítico? ¿Por qué?
\end{itemize}
\end{frame}

\begin{frame}[allowframebreaks,fragile]{Ejercicios Prácticos y Actividades Interactivas}


\textbf{Ejercicio 2: Diseño de Intervención Local (10 minutos)}
\begin{itemize}
  \item \textbf{Basándose en el Ejercicio 1,} su grupo deberá proponer una intervención de fisioterapia en la APS para abordar el factor más crítico identificado.
  \item \textbf{Describan la intervención:}
  \begin{itemize}
    \item ¿En qué consiste? (Ej: programa de ejercicios comunitarios, consulta de fisioterapia a domicilio, taller educativo).
    \item ¿Cómo se implementaría a nivel local?
    \item ¿Qué recursos se necesitarían?
  \end{itemize}
  \item \textbf{Reflexión:} ¿Cómo medirían el éxito o el impacto de esta intervención en la comunidad?
\end{itemize}
\end{frame}

\begin{frame}[allowframebreaks,fragile]{Sabías que... Datos Curiosos y Reveladores}


\begin{itemize}
  \item \textbf{Potencial Subutilizado del Análisis Espacial:} A pesar de su probada eficacia para identificar ``puntos calientes'' de enfermedad (ej. dengue) y correlacionarlos con servicios de APS, su aplicación sistemática para informar la distribución de servicios de fisioterapia sigue siendo limitada. ¡Representa una gran oportunidad!
  \item \textbf{La Fisioterapia como Eje de Equidad en Salud:} El análisis territorial revela cómo una fisioterapia accesible y culturalmente competente en la APS es un pilar fundamental para reducir las desigualdades. Su rol va más allá del tratamiento, abarcando prevención de discapacidad y promoción del envejecimiento activo.
  \item \textbf{Disparidad entre Rol Definido y Práctica Real:} Aunque la literatura destaca el rol multifacético y preventivo del fisioterapeuta en la APS, su integración plena en la práctica diaria a menudo enfrenta barreras. Ejemplos incluyen la falta de reconocimiento institucional, recursos insuficientes o modelos de referencia inadecuados en diferentes territorios.
\end{itemize}
\end{frame}

\begin{frame}[allowframebreaks,fragile]{Conclusiones y Puntos Clave (Takeaways)}


\textbf{El Análisis Territorial: Una Herramienta Indispensable}
\begin{itemize}
  \item Permite una \textbf{visión holística} de las necesidades de la población y las capacidades del sistema de salud en áreas geográficas específicas.
  \item \textbf{Fortalece la toma de decisiones} en salud pública, optimizando la asignación de recursos y el diseño de programas.
\end{itemize}
\textbf{La Fisioterapia en APS: Clave para la Equidad}
\begin{itemize}
  \item Un análisis territorial adecuado visibiliza las necesidades de \textbf{poblaciones vulnerables}, garantizando que la fisioterapia llegue a quienes más la necesitan.
  \item Es un pilar fundamental para la \textbf{prevención, promoción de la salud y el manejo} de condiciones crónicas.
\end{itemize}
\textbf{Cerrando la Brecha}
\begin{itemize}
  \item Es crucial \textbf{cerrar la brecha} entre el potencial teórico y la práctica real de la fisioterapia en la APS.
  \item Esto requiere \textbf{políticas de salud pública} que promuevan la integración, la capacitación en análisis territorial y la asignación de recursos suficientes.
\end{itemize}
\end{frame}

\begin{frame}[allowframebreaks,fragile]{Referencias Bibliográficas}


\begin{itemize}
  \item Ramírez Rosario, M. (2024). P-3. ANÁLISIS DE LA SITUACIÓN ACTUAL DE LOS FISIOTERAPEUTAS DE ATENCIÓN PRIMARIA DEL SERVICIO EXTREMEÑO DE SALUD. \textit{Fisioterapia}. DOI: 10.1016/s0211-5638(24)60063-5
  \item Nascimento, M., \& Rodrigues-Júnior, A. (2014). Análisis espacial del dengue y la atención primaria de salud en Alfenas, Minas Gerais, Brasil. \textit{Atención Primaria, 46}(2), 79-84. DOI: 10.1016/j.aprim.2013.12.016
  \item Zamarreño, F. (2025). Salud, territorio y mujeres trans: miradas desde la Atención Primaria. \textit{ConCienciaSocial}. DOI: 10.66430/3ps38592
  \item Tepox Bruno, R., \& Díaz López, B. (2023). Rol del fisioterapeuta en la atención primaria en salud. \textit{Revista Mexicana de Fisioterapia, 3}(1), 101. DOI: 10.67035/revmexfisiot.vi3.101
\end{itemize}
\end{frame}




\end{document}